In this section, we will prove Theorem \ref{thm:mainboundary}.  Our methods are due to Stephan Stolz \cite{StolzBook}, and we rely heavily on his work.  We are able to improve on Stolz's results by a combination of Theorem \ref{thm:intro-main} and the following:

\begin{thm} \label{thm:mod8}
Suppose that a class $\alpha$ in the $(2k+d)$th homotopy group of the mod $8$ Moore spectrum
    $$\alpha \in \pi_{2k+d}(C(8))$$
has $\mathrm{H}\mathbb{F}_2$-Adams filtration at least $\frac{2k+d}{5}+15$.  Then, if $2k+d \ge 126$, the image of $\alpha$ under the Bockstein map
$$\pi_{2k+d}(C(8)) \to \pi_{2k+d-1}(\mathbb{S})$$
is contained in the subgroup of $\pi_{2k+d-1}(\mathbb{S})$ generated by $\mathcal{J}_{2k+d-1}$ and Adams's $\mu$-family.
\end{thm}

Theorem \ref{thm:mod8} will be proved as Theorem \ref{thm:mod8-main-thm} in the subsequent half of the paper.

\begin{rmk}
    Stolz relied on a similar result for the mod $2$ Moore spectrum, which he attributes to Mahowald \cite[Satz 12.9]{StolzBook}.  Though Mahowald announced such a result in \cite{MahBull}, and again in \cite{MahTokyo} (the reference that Stolz cites), to the best of our knowledge no proof has appeared in print.  Part of our motivation in proving Theorem \ref{thm:mod8} is to fill this gap in the literature. The other motivation is that we obtain stronger geometric consequences from the mod $8$ Moore spectrum.
\end{rmk}

Recall that $\mathrm{MO} \langle k \rangle$ denotes the Thom spectrum of the map
$$\tau_{\ge k} \mathrm{BO} \to \mathrm{BO}.$$
There is a unit map $\mathbb{S} \to \mathrm{MO} \langle k \rangle$, which may be extended to a cofiber sequence
$$\mathbb{S} \to \mathrm{MO} \langle k \rangle \to \mathrm{MO} \langle k \rangle / \mathbb{S} \stackrel{\partial}{\to} \mathbb{S}^1.$$
Stolz constructed \cite[Satz 3.1]{StolzBook} a spectrum $A[k]$ together with a map $$b:A[k] \to \mathrm{MO} \langle k \rangle / \mathbb{S}$$ such that the following is true:

\begin{thm}[{\cite[Lemma 12.5]{StolzBook}}]
Let $k>2$ and $d \ge 0$ be integers.  Suppose that, for every element $\alpha \in \pi_{2k+d}(A[k])$, the image of $\alpha$ under the composite
$$\pi_{2k+d}(A[k]) \stackrel{b_*}{\longrightarrow} \pi_{2k+d} \left( \mathrm{MO} \langle k \rangle / \mathbb{S} \right) \stackrel{\partial_*}{\longrightarrow} \pi_{2k+d}(\mathbb{S}^1) \cong \pi_{2k+d-1}(\mathbb{S})$$
is in $\mathcal{J}_{2k+d-1}$.  Then the boundary of any $(k-1)$-connected, almost closed $(2k+d)$-manifold also bounds a parallelizable manifold.
\end{thm}

Stolz then proved the following two theorems. The first is contained in the proof of \cite[Satz 12.7]{StolzBook}:

\begin{thm}\cite[Proof of Satz 12.7]{StolzBook}  \label{thm:Stolz-Adams}
    Suppose that $k>2$ and $d \ge 0$, and let $M_{2k+d+1} \hookrightarrow A[k]$ denote a $(2k+d+1)$-skeleton of $A[k]$. Then the composite \[M_{2k+d+1} \hookrightarrow A[k] \to \mathrm{MO} \langle k \rangle / \Ss \to \Ss^1\] has $\mathrm{H}\mathbb{F}_2$-Adams filtration at least $$h(k-1)-\lfloor \mathrm{log}_2(2k+d+1) \rfloor +1,$$ where $h(k-1)$ is the number of integers of $s$ with $0<s \le k-1$ and $s \equiv 0,1,2,$ or $4$ modulo $8$.
\end{thm}

The second follows from \cite[Theorem A]{StolzBook} along with the equivalences $A[k] \simeq A[k+1]$ for $k \equiv 3,5,6,7 \mod 8$ \cite[Satz 3.1(ii)]{StolzBook}:

\begin{thm}[{\cite[Theorem A]{StolzBook}}] \label{thm:Stolz-8torsion}
Suppose that $k \ge 9$ and that $0 \le d \le 3$.  Then, unless one of the following conditions are met, every element of $\pi_{2k+d} A[k]$ is $8$-torsion:
\begin{itemize}
\item $k \equiv 0$ modulo $4$ and $d=0$.
\item $k \equiv 3$ modulo $4$ and $d=2$.
\end{itemize}
\end{thm}

We now proceed with our proof of Theorem \ref{thm:mainboundary}.

\begin{thm} \label{thm:strongmainboundary}
Let $k>124$ and $0 \le d \le 3$ be integers.  Suppose that $k \equiv 0$ modulo $4$ and $d=0$ or $k \equiv 3$ modulo $3$ and $d=2$.  Then the boundary of any $(k-1)$-connected, almost closed $(2k+d)$-manifold also bounds a parallelizable manifold.
\end{thm}

\begin{proof}
We need to show that the image of the composite
$$\pi_{2k+d} A[k] \to \pi_{2k+d}(\MO \langle k \rangle/\mathbb{S}) \to \pi_{2k+d}(\mathbb{S}^1)$$
contains only classes in $\mathcal{J}_{2k+d-1}$.  In fact, we will show that this is already true of the image of
$$\pi_{2k+d}(\MO \langle k \rangle/\mathbb{S}) \to \pi_{2k+d}(\mathbb{S}^1),$$
or equivalently that the unit map
$$\pi_{2k+d-1} \mathbb{S} \to \pi_{2k+d-1} \MO \langle k \rangle$$
has kernel consisting only of classes in $\mathcal{J}_{2k+d-1}$.  If $d=0$ and $k \equiv 0$ modulo $4$, this follows from Theorem \ref{thm:intro-main}.  If $d=2$ and $k \equiv 3$ modulo $4$, then $\MO \langle k \rangle \simeq \MO \langle k+1 \rangle$, and so $\pi_{2k+2} \MO \langle k \rangle \cong \pi_{2k+2} \MO \langle k+1 \rangle$ and this again follows from Theorem \ref{thm:intro-main}.
\end{proof}

\begin{thm} \label{thm:weakmainboundary}
Let $k>232$ and $0 \le d \le 3$ be integers.  Suppose that $k$ and $d$ satisfy neither of the exceptional conditions under which Theorem \ref{thm:Stolz-8torsion} fails and Theorem \ref{thm:strongmainboundary} applies.  Then the boundary of any $(k-1)$-connected, almost closed $(2k+d)$-manifold also bounds a parallelizable manifold. 
\end{thm}

\begin{proof}
We construct an argument very similar to that found on \cite[p.107]{StolzBook}.  Namely, consider the diagram
    \[
\begin{tikzcd}
    \Sigma^{-1} C(8) \otimes M_{2k+d+1} \arrow{r}{1 \otimes \iota} \arrow{d} & \Sigma^{-1}C(8) \otimes A[k] \arrow{r}{1 \otimes (\partial \circ b)} \arrow{d} & \Sigma^{-1}C(8) \otimes \mathbb{S}^1 \arrow{d} \\   
    M_{2k+d+1} \arrow{r}{\iota} & A[k] \arrow{r}{\partial \circ b} &  \mathbb{S}^1
\end{tikzcd}
    \]
    where $M_{2k+d+1} \to A[k]$ is a $(2k+d+1)$-skeleton.

    Let $\alpha$ denote a map $\Ss^{2k+d} \to A[k]$. Then we may factor $\alpha$ through an $8$-torsion map $\Ss^{2k+d} \to M_{2k+d+1}$ by \Cref{thm:Stolz-8torsion}, and thus we may choose a lift $\bar{\alpha}:\Ss^{2k+d} \to \Sigma^{-1}C(8) \otimes M_{2k+d+1}$. Since \[M_{2k+d+1} \xrightarrow{\iota} A[k] \xrightarrow{\partial\circ b} \Ss^{1}\] is of $\HFt$-Adams filtration at least $h(k-1)-\lfloor \mathrm{log}_2(2k+d+1) \rfloor +1$ by \Cref{thm:Stolz-Adams}, so is \[\Sigma^{-1}C(8) \otimes M_{2k+d+1} \xrightarrow{1 \otimes \iota} \Sigma^{-1} C(8) \otimes A[k] \xrightarrow{1 \otimes (\partial \circ b)} \Sigma^{-1} C(8) \otimes \mathbb{S}^1.\] It follows that $(1 \otimes \partial) \circ (1 \otimes b) \circ (1 \otimes \iota) \circ \bar{\alpha} \in \pi_{2k+d} \left(C(8)\right)$ is too. Thus, by Theorem \ref{thm:mod8}, so long as
$$2k+d \ge 126, \text{ and}$$
$$h(k-1)-\lfloor \mathrm{log}_2(2k+d+1) \rfloor +1 \ge \frac{2k+d}{5}+15,$$
the image of $\alpha$ in $\pi_{2k+d-1} \mathbb{S}$ must be in the subgroup generated by $\mathcal{J}_{2k+d-1}$ and Adams's $\mu$-family.  In \Cref{lemm:inequality}, we show that both of these conditions are satisfied under our assumptions $k > 232$ and $0 \le d \le 3$.
Now, we claim that the image of $\alpha$ in $\pi_{2k+d-1} \mathbb{S}$ must actually be in the subgroup generated by $\mathcal{J}_{2k+d-1}$, without Adams's $\mu$-family.  This follows simply from the fact that this class is in the image of the map
$$\pi_{2k+d} \MO \langle k \rangle / \mathbb{S} \to \pi_{2k+d} \mathbb{S}^1,$$
and therefore in the kernel of the map
$$\pi_{2k+d-1} \mathbb{S} \to \MO \langle k \rangle.$$
Recall that the Atiyah--Bott--Shapiro orientation \cite{ABS} determines a unital map 
$$\mathrm{MO} \langle 3 \rangle = \mathrm{MSpin} \to \mathrm{KO}.$$
The composite map
$$\pi_{2k+d-1} \mathbb{S} \to \pi_{2k+d-1} \MO \langle k \rangle \to \pi_{2k+d-1} \MO \langle 3 \rangle \to \pi_{2k+d-1} \mathrm{KO}$$
has the effect of killing $\mathcal{J}_{2k+d-1}$ without killing any of Adams's $\mu$-family.  Thus, any class in the kernel of the map
$$\pi_{2k+d-1} \mathbb{S} \to \pi_{2k+d-1} \MO \langle k \rangle$$
cannot be a sum of a class in $\mathcal{J}_{2k+d-1}$ and a non-trivial element of the $\mu$-family.
\end{proof}

\begin{proof} [Proof of Theorem \ref{thm:mainboundary}]
This follows by combining Theorem \ref{thm:strongmainboundary} and Theorem \ref{thm:weakmainboundary}.
\end{proof}

\begin{lem}\label{lemm:inequality}
    The inequality \[h(k-1) - \lfloor \log_2(2k+d+1) \rfloor + 1 \geq \frac{2k+d}{5} + 15\] holds for $k > 232$ and $0 \leq d \leq 3$.
\end{lem}

\begin{proof}
  Without loss of generality we may assume $d=3$.
  Then, using the inequality $h(k-1) \geq \frac{k}{2}-1$, it will suffice to show that
  % \[\frac{k}{2} \geq \frac{2k+d}{5} + 15 + \log_2 (2k+d+1),\] which by using $0 \leq d \leq 3$ we may reduce to
  \begin{align}
    \frac{k}{10} \geq \frac{3}{5} + 15 + \log_2 (2k+4). \label{eq:stolz-inequality}
  \end{align}
  Taking derivatives we see that the left hand side increases faster than the right hand side as soon as $k \geq 13$. Using a computer we find that \Cref{eq:stolz-inequality} holds for $k=246$, so the lemma holds for $k \geq 246$. For the remaining values of $k$, we compute each side of the desired inequality
  % Using a computer, we find that for $k = 246$ we have
  % \[\frac{246}{10} = 24.6 \geq \frac{3}{5} + 15 + \log_2 (496) \approx 23.5542\]
  % so the lemma holds for $k \geq 246$.
  \[h(k-1) - \lfloor \log_2(2k+d+1) \rfloor + 1 \geq \frac{2k+d}{5} + 15\]
  for $d=3$, and display their difference, $\Delta$, in the following table:

  \begin{table}[ht!]
    \begin{center}
      \begin{tabular}{|c||c|c|c|c|c|c|c|c|c|c|c|c|c|}\hline
        $k$ & 233 & 234 & 235 & 236 & 237 & 238 & 239 & 240 & 241 & 242 & 243 & 244 & 245 \\\hline
        $\Delta$ & 0.2 & 0.8 & 1.4 & 1.0 & 1.6 & 1.2 & 0.8 & 0.4 & 1.0 & 1.6 & 2.2 & 1.8 & 2.4 \\\hline
      \end{tabular}
    \end{center}
  \end{table}
  
%% \begin{table}[h!]
%% \begin{center}
%% \begin{tabular}{|c||c|c|}\hline
%% $k$ & $h(k-1) - \lfloor \log_2(2k+4) \rfloor + 1$ & $\frac{2k+3}{5} + 15$ \\
%% \hline\hline
%% 233 & 109 & 108.8 \\\hline
%% 234 & 110 & 109.2 \\\hline
%% 235 & 111 & 109.6 \\\hline
%% 236 & 111 & 110 \\\hline
%% 237 & 112 & 110.4 \\\hline
%% 238 & 112 & 110.8 \\\hline
%% 239 & 112 & 111.2 \\\hline
%% 240 & 112 & 111.6 \\\hline
%% 241 & 113 & 112 \\\hline
%% 242 & 114 & 112.4 \\\hline
%% 243 & 115 & 112.8 \\\hline
%% 244 & 115 & 113.2 \\\hline
%% 245 & 116 & 113.6 \\\hline
%% \end{tabular}
%% \end{center}
%% \end{table}

\end{proof}

\begin{rmk} \label{rmk:open-questions}
In Section \ref{sec:open-problems} we discussed possible improvements to Theorem \ref{thm:strongmainboundary}.
We have spent comparatively little effort optimizing Theorem \ref{thm:weakmainboundary}, and it would be interesting to see an improvement of the bounds $k>232$ and $d \le 3$.

For a fixed dimension $m$, it would be interesting to know the largest integer $\ell$ such that a smooth, $\ell$-connected, almost closed $m$-manifold bounds an element non-trivial in $\mathrm{coker}(J)_{m-1}$.
Conjecture \ref{cnj:vanishing-line} suggests that this integer $\ell$ should be closer to $\frac{m}{3}$ than $\frac{m}{2}$.
\end{rmk}

%%% Local Variables:
%%% mode: latex
%%% TeX-master: "Main"
%%% eval: (visual-line-mode t)
%%% eval: (auto-fill-mode 0)
%%% End:
